% ============================================================
%  MTH636 – Game Theory
%  Short Notes / Exam Prep Guide
% ============================================================
\documentclass[12pt, a4paper]{article}

% ============================================================
%  MTH636 -- Game Theory
%  Notes Preamble — Color-Coded, TikZ-Enhanced
% ============================================================

% --- Core packages ---
\usepackage[T1]{fontenc}
\usepackage[utf8]{inputenc}
\usepackage{lmodern}
\usepackage[margin=2.5cm, top=2.8cm, bottom=2.8cm]{geometry}
\usepackage{microtype}

% --- Math ---
\usepackage{amsmath, amssymb, amsthm, mathtools}
\usepackage{bm}        % bold math
\usepackage{cancel}    % strike-through in math

% --- Color & Boxes ---
\usepackage[dvipsnames, svgnames, x11names, table]{xcolor}
\usepackage[most]{tcolorbox}
\tcbuselibrary{breakable, skins, theorems}

% --- TikZ & Diagrams ---
\usepackage{tikz}
\usepackage{tikz-qtree}
\usetikzlibrary{arrows.meta, positioning, shapes.geometric, fit, backgrounds,
                decorations.pathreplacing, calc, matrix, trees}
\usepackage{pgfplots}
\pgfplotsset{compat=1.18}

% --- Tables ---
\usepackage{booktabs}
\usepackage{multirow}
\usepackage{array}
\usepackage{tabularx}
\usepackage{makecell}

% --- Typography & Layout ---
\usepackage{enumitem}
\usepackage{parskip}
\usepackage{fancyhdr}
\usepackage{titlesec}
\usepackage{titletoc}
\usepackage{hyperref}
\usepackage{graphicx}
\usepackage{float}
\usepackage{caption}

% --- Color Palette ---
\definecolor{defblue}{HTML}{1A5276}      % Definition header
\definecolor{defbluebg}{HTML}{D6EAF8}    % Definition background
\definecolor{thmgreen}{HTML}{1E8449}     % Theorem header
\definecolor{thmbg}{HTML}{D5F5E3}        % Theorem background
\definecolor{exorange}{HTML}{B7770D}     % Example header
\definecolor{exbg}{HTML}{FEF9E7}         % Example background
\definecolor{warnred}{HTML}{922B21}      % Warning/Key Insight header
\definecolor{warnbg}{HTML}{FADBD8}       % Warning background
\definecolor{noticegray}{HTML}{2C3E50}   % Note header
\definecolor{noticebg}{HTML}{EBF5FB}     % Note background
\definecolor{taketawayyellow}{HTML}{7D6608} % Key Takeaway header
\definecolor{takeawaybg}{HTML}{FDFEFE}   % Takeaway background
\definecolor{accentpurple}{HTML}{6C3483} % Accent / proof
\definecolor{lightgray}{HTML}{F2F3F4}    % shading

% ============================================================
%  TCOLORBOX ENVIRONMENTS
% ============================================================

% -- Definition (Blue) --
\newtcolorbox{defbox}[2][]{
  enhanced, breakable,
  colback=defbluebg, colframe=defblue, coltitle=white,
  fonttitle=\bfseries\sffamily,
  title={Definition: #2},
  attach boxed title to top left={yshift=-2mm, xshift=4mm},
  boxed title style={colback=defblue, rounded corners},
  arc=4pt, boxrule=1.2pt,
  left=6pt, right=6pt, top=8pt, bottom=6pt,
  #1
}

% -- Theorem / Proposition (Green) --
\newtcolorbox{thmbox}[2][]{
  enhanced, breakable,
  colback=thmbg, colframe=thmgreen, coltitle=white,
  fonttitle=\bfseries\sffamily,
  title={Theorem: #2},
  attach boxed title to top left={yshift=-2mm, xshift=4mm},
  boxed title style={colback=thmgreen, rounded corners},
  arc=4pt, boxrule=1.2pt,
  left=6pt, right=6pt, top=8pt, bottom=6pt,
  #1
}

% -- Proposition (Green variant) --
\newtcolorbox{propbox}[2][]{
  enhanced, breakable,
  colback=thmbg, colframe=thmgreen!70!black, coltitle=white,
  fonttitle=\bfseries\sffamily,
  title={Proposition: #2},
  attach boxed title to top left={yshift=-2mm, xshift=4mm},
  boxed title style={colback=thmgreen!70!black, rounded corners},
  arc=4pt, boxrule=1.2pt,
  left=6pt, right=6pt, top=8pt, bottom=6pt,
  #1
}

% -- Example (Orange/Yellow) --
\newtcolorbox{exbox}[2][]{
  enhanced, breakable,
  colback=exbg, colframe=exorange, coltitle=white,
  fonttitle=\bfseries\sffamily,
  title={Example: #2},
  attach boxed title to top left={yshift=-2mm, xshift=4mm},
  boxed title style={colback=exorange, rounded corners},
  arc=4pt, boxrule=1.2pt,
  left=6pt, right=6pt, top=8pt, bottom=6pt,
  #1
}

% -- Key Insight / Warning (Red) --
\newtcolorbox{warnbox}[2][]{
  enhanced, breakable,
  colback=warnbg, colframe=warnred, coltitle=white,
  fonttitle=\bfseries\sffamily,
  title={Key Insight: #2},
  attach boxed title to top left={yshift=-2mm, xshift=4mm},
  boxed title style={colback=warnred, rounded corners},
  arc=4pt, boxrule=1.2pt,
  left=6pt, right=6pt, top=8pt, bottom=6pt,
  #1
}

% -- Note (Gray-Blue) --
\newtcolorbox{notebox}[1][]{
  enhanced, breakable,
  colback=noticebg, colframe=noticegray!60, coltitle=white,
  fonttitle=\bfseries\sffamily,
  title={Note},
  attach boxed title to top left={yshift=-2mm, xshift=4mm},
  boxed title style={colback=noticegray!60, rounded corners},
  arc=4pt, boxrule=0.8pt,
  left=6pt, right=6pt, top=8pt, bottom=6pt,
  #1
}

% -- Key Takeaway --
\newtcolorbox{takebox}[1][]{
  enhanced, breakable,
  colback=takeawaybg, colframe=taketawayyellow, coltitle=white,
  fonttitle=\bfseries\sffamily,
  title={Key Takeaway},
  attach boxed title to top left={yshift=-2mm, xshift=4mm},
  boxed title style={colback=taketawayyellow, rounded corners},
  arc=4pt, boxrule=1.2pt,
  left=6pt, right=6pt, top=8pt, bottom=6pt,
  #1
}

% -- Intuition --
\newtcolorbox{intubox}[1][]{
  enhanced, breakable,
  colback=accentpurple!8, colframe=accentpurple!60, coltitle=white,
  fonttitle=\bfseries\sffamily,
  title={Intuition},
  attach boxed title to top left={yshift=-2mm, xshift=4mm},
  boxed title style={colback=accentpurple!80, rounded corners},
  arc=4pt, boxrule=0.8pt,
  left=6pt, right=6pt, top=8pt, bottom=6pt,
  #1
}

% ============================================================
%  PAYOFF MATRIX HELPER
%  Usage: inside tabular, use \payoff{a}{b} for (a,b)
% ============================================================
\newcommand{\payoff}[2]{$(#1,\;#2)$}
\newcommand{\nash}[2]{\underline{$#1$},\underline{$#2$}}  % underline NE payoffs
\newcommand{\Nset}{\mathbb{N}}
\newcommand{\Rset}{\mathbb{R}}
\newcommand{\Sset}{\mathcal{S}}
\newcommand{\Ui}{u_i}

% ============================================================
%  SECTION FORMATTING
% ============================================================
\titleformat{\section}
  {\large\bfseries\sffamily\color{defblue}}
  {\thesection}{1em}{}[\titlerule]

\titleformat{\subsection}
  {\normalsize\bfseries\sffamily\color{thmgreen!80!black}}
  {\thesubsection}{1em}{}

\titleformat{\subsubsection}
  {\normalsize\itshape\sffamily\color{exorange!80!black}}
  {\thesubsubsection}{1em}{}

% ============================================================
%  HEADER / FOOTER
% ============================================================
\pagestyle{fancy}
\fancyhf{}
\fancyhead[L]{\sffamily\small\color{defblue} MTH636 -- Game Theory}
\fancyhead[R]{\sffamily\small\color{defblue} \leftmark}
\fancyfoot[C]{\sffamily\small\thepage}
\renewcommand{\headrulewidth}{0.4pt}

% ============================================================
%  HYPERREF SETUP
% ============================================================
\hypersetup{
  colorlinks=true,
  linkcolor=defblue,
  urlcolor=thmgreen,
  citecolor=accentpurple,
  bookmarksnumbered=true,
  pdftitle={MTH636 Game Theory Notes},
  pdfauthor={MTH636 Student Notes}
}

% ============================================================
%  MISC HELPERS
% ============================================================
\newcommand{\BR}{\text{BR}}        % Best Response
\newcommand{\NE}{\text{NE}}        % Nash Equilibrium
\newcommand{\SPNE}{\text{SPNE}}    % Subgame Perfect NE
\newcommand{\Dom}{\text{dom}}      % Dominance
\newcommand{\mixs}[1]{\sigma_{#1}} % Mixed strategy
\newcommand{\pures}[1]{s_{#1}}     % Pure strategy


% ============================================================
\begin{document}

% Title
\begin{center}
  {\LARGE\bfseries\sffamily\color{defblue} Game Theory: Short Notes}\\[6pt]
  {\Large\sffamily Intuition, Frameworks, and Problem-Solving Guide}\\[4pt]
\end{center}
\vspace{0.3cm}
\hrule height 1.5pt
\vspace{0.6cm}

% ============================================================
\section{The Core Philosophy of Game Theory}
Game theory isn't just about finding the ``best'' move in a vacuum; it is about \textbf{interactive decision making}. 
\textit{Core Intuition:} You must put yourself in your opponent's shoes. ``What will they do? If they do X, what is my best response? But they know I will play my best response to X, so will they actually play X?''

% ============================================================
\section{Normal Form (Simultaneous) Games}

A game needs 3 ingredients: Players ($N$), Strategies ($S_i$), and Payoff/Utility functions ($u_i$).

\subsection{Dominance and IESDS}
\begin{itemize}
    \item \textbf{Strictly Dominated Strategy:} A strategy $A$ is strictly dominated by $B$ if playing $B$ \textit{always} gives you a strictly higher payoff than playing $A$, regardless of what your opponent does.
    \item \textbf{Rationality Assumption:} A rational player will \textbf{never} play a strictly dominated strategy.
    \item \textbf{IESDS (Iterated Elimination of Strictly Dominated Strategies):} 
    \textit{How to use it:} Cross out strictly dominated rows for P1. Then, in the remaining matrix, cross out strictly dominated columns for P2. Repeat until nothing else can be crossed out. 
    \textit{Trick:} If an exam question asks you to solve a large $4 \times 4$ game, always check for dominated strategies first. It usually collapses into a $2 \times 2$ game!
    \item \textbf{Domination by a Mixed Strategy:} A pure strategy can be strictly dominated by a \textit{mixture} of other strategies! (e.g., $1/2(\text{Top}) + 1/2(\text{Bottom})$ strictly dominates Middle). If it is strictly dominated (even by a mixed strategy), it will \textbf{never} get positive probability in any Nash Equilibrium. Always check for mixed domination if you can't eliminate a strategy purely.
\end{itemize}

\subsection{Pure Strategy Nash Equilibrium (NE)}
\textbf{Intuition:} A state of the world where \textbf{no one has regrets}. If everyone's choice is revealed, no single player would want to change their move.
\begin{itemize}
    \item \textbf{Best Response (BR):} The strategy that yields the highest payoff given a specific belief about what the opponent will do.
    \item \textbf{How to find Pure NE (The Underline Method):}
    \begin{enumerate}
        \item Assume P2 plays Left. Underline P1's highest payoff in that column. Repeat for all columns.
        \item Assume P1 plays Top. Underline P2's highest payoff in that row. Repeat for all rows.
        \item Any cell with \textbf{both payoffs underlined} is a Nash Equilibrium!
    \end{enumerate}
\end{itemize}

\begin{notebox}
  \textbf{Classic Game Archetypes (Learn the stories, not just the matrices):}
  \begin{itemize}
      \item \textbf{Prisoner's Dilemma:} Individual rationality leads to collective ruin. (Defect, Defect) is the unique NE, but it is Pareto inefficient.
      \item \textbf{Battle of the Sexes:} Players want to coordinate, but disagree on where. Yields multiple Pure NEs (coordination points) and one Mixed NE.
      \item \textbf{Matching Pennies:} Pure conflict (Zero-sum). P1 wants to match, P2 wants to mismatch. Result: \textbf{No Pure NE exists}. You must be unpredictable (mix).
  \end{itemize}
\end{notebox}

% ============================================================
\section{Mixed Strategies \& The Indifference Principle}

\begin{thmbox}{Nash's Existence Theorem}
  \textbf{Every finite game has at least one Nash Equilibrium.} If you can't find a Pure NE, a Mixed NE is mathematically guaranteed to exist (proved via Brouwer's Fixed Point Theorem).
\end{thmbox}

If a game has no Pure NE (or if you need to find all NEs), players randomize their choices using probabilities.

\textbf{The Indifference Principle (The Golden Rule of Mixing):} \\
If a player is randomizing between two strategies (say, Top and Bottom), it \textit{must} be because both strategies give them the \textbf{exact same expected payoff} against the opponent's mixed strategy. If one gave a higher expected payoff, the player would just play that one 100\% of the time!

\begin{thmbox}{How to Solve for Mixed NE (Crucial Step-by-Step)}
  Let P1 play Top with probability $p$ and Bottom with $1-p$. \\
  Let P2 play Left with probability $q$ and Right with $1-q$.
  \begin{enumerate}
      \item \textbf{To find P2's mixing prob ($q$):} Set P1's expected payoffs equal!
      \[ E[U_1(\text{Top})] = E[U_1(\text{Bottom})] \]
      Express these expected utilities in terms of $q$. Solve for $q$.
      \item \textbf{To find P1's mixing prob ($p$):} Set P2's expected payoffs equal!
      \[ E[U_2(\text{Left})] = E[U_2(\text{Right})] \]
      Express these in terms of $p$. Solve for $p$.
  \end{enumerate}
\end{thmbox}
\textit{Deep Intuition:} Your mixing probabilities are designed to keep your opponent indifferent, \textbf{not} to maximize your own payoff in isolation! You mix to make the opponent indifferent so they don't exploit you.

% ============================================================
\section{Extensive Form (Sequential) Games}

When players move one after another, we use Game Trees.
\begin{itemize}
    \item \textbf{Nodes \& Branches:} Decision points and actions.
    \item \textbf{Information Sets:} A dotted line (or bubble) connecting nodes. It means the player whose turn it is \textit{cannot distinguish} which of those nodes they are actually at. (e.g., they don't know what the previous player chose).
    \item \textbf{Perfect Information:} Every information set contains exactly one node. (Like Chess).
\end{itemize}

\subsection{Backward Induction and SPNE}
If you just convert a tree into a matrix and find the Nash Equilibria, you will often find NEs that rely on \textbf{Non-Credible Threats}. (e.g., ``If you enter the market, I will start a price war.'' But if the price war hurts both, would they actually do it if tested?)

\textbf{Subgame Perfect Nash Equilibrium (SPNE):} \\
A refinement of NE that eliminates non-credible threats. It requires that the strategies constitute a Nash Equilibrium in \textit{every subgame} (every smaller subtree).

\begin{thmbox}{How to Solve for SPNE (Backward Induction)}
  \begin{enumerate}
      \item Go to the very end of the tree (the final decision nodes before the payoffs).
      \item Ask: ``What is the rational choice for the player at this node?'' Highlight that branch.
      \item Replace that entire subgame with the payoff of the highlighted branch.
      \item Move one step up the tree and repeat.
  \end{enumerate}
  \textit{Note:} Backward induction only works perfectly for games of \textit{perfect information}. If there are information bubbles, you must treat that bubble as a simultaneous mini-game (matrix) and find its NE, then fold that NE payoff back up the tree.
\end{thmbox}

% ============================================================
\section{Emergency Exam Tips}
\begin{itemize}
    \item \textbf{If asked to find ALL Nash Equilibria:} Don't forget the mixed ones! Check pure NE first, then assign probabilities, apply the indifference principle, and check for mixed NE.
    \item \textbf{If a mixed probability equation gives $q = 2$ or $q < 0$:} You've made an arithmetic mistake, or there is no mixed NE involving those specific strategies. Probabilities must be between 0 and 1.
    \item \textbf{If a game looks too big to solve:} Look immediately for strictly dominated strategies to delete. 
    \item \textbf{If asked why a threat is ``non-credible'':} Show that if the subgame was actually reached, the player would strictly prefer to take a different action than the one they threatened.
\end{itemize}

\end{document}
